%% kapa: exam=1 round=31 year=2020 period=2 form=A subject=accounting subject_ko=회계학 seq=12 number=52
%% source: Q-Net 감정평가사 기출문제 (cst00309) artlSeq=5207219 file=extracted/1008/5207219/A형/A형/2020년도 제31회 감정평가사 1차 시험_2교시_A형.pdf
%% conversion: kordoc 4.24.51 -> tools/gen1.py (auto); review=proofread-sonnet; flags=none
\begin{question}{52}
(주)감평은 확정급여제도를 채택하고 있으며, 20x1년 초 순확정급여부채는 ￦20,000이다. (주)감평의 20x1년도 확정급여제도와 관련된 자료는 다음과 같다.
\begin{framedtext}
○ 순확정급여부채(자산) 계산시 적용한 할인율은 연 6\%이다. \par
○ 20x1년도 당기근무원가는 ￦85,000이고, 20x1년 말 퇴직종업원에게 ￦38,000의 현금이 사외적립자산에서 지급되었다. \par
○ 20x1년 말 사외적립자산에 ￦60,000을 현금으로 출연하였다. \par
○ 20x1년에 발생한 확정급여채무의 재측정요소(손실)는 ￦5,000이고, 사외적립자산의 재측정요소(이익)는 ￦2,200이다. \par
\end{framedtext}
(주)감평이 20x1년 말 재무상태표에 순확정급여부채로 인식할 금액과 20x1년도 포괄손익계산서상 당기손익으로 인식할 퇴직급여 관련 비용은?
{\par\noindent\hspace*{2em}\makebox[4cm][c]{\uline{순확정급여부채}}\makebox[4cm][c]{\uline{퇴직급여 관련 비용}}\par}
\begin{choices}
  \item \makebox[4cm][c]{￦11,000}\makebox[4cm][c]{￦85,000}
  \item \makebox[4cm][c]{￦11,000}\makebox[4cm][c]{￦86,200}
  \item \makebox[4cm][c]{￦43,400}\makebox[4cm][c]{￦86,200}
  \item \makebox[4cm][c]{￦49,000}\makebox[4cm][c]{￦85,000}
  \item \makebox[4cm][c]{￦49,000}\makebox[4cm][c]{￦86,200}
\end{choices}
\end{question}
\answer{5}
